% Abstract — structured summary of the monograph
\begin{abstract}
\noindent\textbf{Schema coherence is a measurable property of learned representation
structure whose deterioration can produce a stable low-coherence regime (the
$\sigma$-trap) associated with systematic generalisation failure.} This monograph
develops and defends that claim in three cumulative tiers.

\textbf{Diagnosis (Ch.~2--3).} A scoping review of the AGI safety landscape and a
systematic review of the $\sigma$-trap evidence base show that the safety literature is
fragmented and that internal representation structure is treated operationally rather
than as an explicit safety property; the reviews map the terrain without presupposing
the thesis.

\textbf{Framework and evidence (Ch.~4--8).} The $\Sigma$-Align framework operationally
defines schema coherence as a measurable, intervenable quantity and submits it to
construct-validity tests. A pilot study tests the causal chain
intervention $\rightarrow \sigma \rightarrow$ generalisation, a meta-analysis
quantifies the $\sigma$-trap within predefined effect families, and two empirical
studies show that a minimal learning system can exhibit a transition from compositional
to non-compositional learning predicted by an independently measurable $\sigma$ (the
$\Sigma$-Model), and that the mechanism has safety-relevant consequences under explicit
diagnostic criteria.

\textbf{Implications (Ch.~9--10).} The evidence justifies a bounded conclusion:
schema-coherent training may be a candidate component of alignment strategy, with
limitations stated before implications.

\textbf{Falsifiability:} the mechanism claim is falsified if $\sigma$ is measured,
interventions raise $\sigma$, and the predicted outcome does not change (e.g.\ a
low-$\sigma$ regime with intact systematic generalisation).

\textbf{Keywords:} compositional generalisation, schema coherence, AI alignment, AGI safety,
neural ODEs, dynamical systems, mesa-optimization, curriculum learning.
\end{abstract}
